\documentclass[12pt]{article}
\usepackage[utf8]{inputenc}
\usepackage[T1]{fontenc}
\usepackage{hyperref}
\title{Wutung Language Resources for K-12 Curriculum Development}
\author{Compiled by Claude Sonnet 5.5 (claude-sonnet-5-5) for Andr\'as Kornai}
\date{}
\begin{document}
\maketitle

\section{Introduction}
Wutung (ISO 639-3 wut, Glottocode wutu1244) is a Skou language of the Western Skou branch. It is spoken on the northern coast of Papua New Guinea, in Sandaun Province near the Indonesian border, with some speakers across the border in the Jayapura area. Wikipedia lists about 900 speakers (2015). Wutung, Musu and Nyao are very closely related, mutually intelligible varieties. Some sources treat Sangke (Nyao) as a separate language. Tok Pisin and English are widely spoken locally. Overviews are at \url{https://en.wikipedia.org/wiki/Wutung_language} and \url{https://en.wikipedia.org/wiki/Skou_languages}.

\section{Primary Educational Resources}
Searches found no primers, textbooks or school materials in Wutung. The language is not known to be taught in schools. A curriculum would have to be built from scratch with the community.

\section{Linguistic and Reference Materials}
Doug Marmion's 2010 ANU PhD thesis, \emph{Topics in the Phonology and Morphology of Wutung}, is the main descriptive source found: \url{https://openresearch-repository.anu.edu.au/items/fc6b56cf-b262-420e-a171-bce6978b6836/full}. All Skou languages have contrastive tone, which any orthography will have to deal with. The Endangered Languages Project has an entry: \url{https://endangeredlanguages.com/lang/10587?hl=es}.

\section{Religious and Community Resources}
The Joshua Project lists the Wutung (Udung) people group: \url{https://joshuaproject.net/index.php/people_groups/15426/PP}. No Wutung scripture or community organisation was found in the search results.

\section{Digital and Audio Resources}
No audio, video, corpus or software resources specific to Wutung were found in these searches.

\section{Curriculum Development Context}
Very little published material exists. The Marmion thesis could support developing an orthography and a word list. Any curriculum would need work with the Wutung community. It would probably use Tok Pisin and English as bridging languages.

\end{document}